\section{ImageNet-9 Protocol Robustness Checks}
\label{app:imagenet9-results}

\begin{figure*}[!t]
\centering
\begin{subfigure}[t]{0.45\textwidth}
    \centering
    \includegraphics[width=\linewidth]{figures/appendix/imagenet9/fig_app_patch_order_stability.pdf}
    \caption{Patch-order stability.}
\end{subfigure}
\hspace{0.06\textwidth}
\begin{subfigure}[t]{0.43\textwidth}
    \centering
    \includegraphics[width=\linewidth]{figures/appendix/imagenet9/fig_app_threshold_sensitivity.pdf}
    \caption{Endpoint-threshold sensitivity.}
\end{subfigure}
\caption{\textbf{Robustness of the ImageNet-9 protocol audit.}
(a) The two independently tie-broken nested-patch orders agree almost perfectly for every response component.
(b) Varying the endpoint threshold changes the allocation between active and endpoint-null response, but not the qualitative protocol conclusion.}
\label{fig:app-in9-sensitivity}
\end{figure*}


The protocol audit in Section~\ref{sec:imagenet9} changes the reveal path while keeping the factual--counterfactual endpoints fixed. We retain two checks that test whether the reported effect is caused by an arbitrary implementation choice.

First, the nested-patch reveal is evaluated with two independently tie-broken patch orders. Their response summaries agree almost perfectly, showing that the blend-to-patch contrast is not an accident of one patch sequence. Second, we recompute the decomposition at $\varepsilon\in\{0.01,0.02,0.05\}$. Changing the threshold reallocates response between the active and endpoint-null branches as expected, but preserves the qualitative conclusion that patch reveal increases total response without increasing evidence.

